\section*{Bab \ref{ch:b3:atmospheric-dynamics} --- Dinamika Atmosfer}

\begin{solution}{exo:b3:atmospheric-dynamics:1}
(a)~$\zeta = 2\omega = 2\times2\pi/86\,400 = \qty{1.45e-4}{s^{-1}}$.
(b)~$f = \num{7.29e-5}$ dan \qty{1.26e-4}{s^{-1}}: pusaran memiliki kira-kira
dua kali \omterm{def:b3:atmospheric-dynamics:vorticity-kinds}{vortisitas planeter} pada \ang{30}~LU dan sedikit lebih besar dari itu pada
\ang{60}~LU.
\end{solution}

\begin{solution}{exo:b3:atmospheric-dynamics:2}
$\zeta = f_{30} - f_{40} = \num{7.29e-5} - \num{9.37e-5} =
\qty{-2.08e-5}{s^{-1}}$: antisiklonik.
\end{solution}

\begin{solution}{exo:b3:atmospheric-dynamics:3}
$k = 2\pi/(5\times10^6) = \qty{1.257e-6}{m^{-1}}$; $c = -\beta/k^2 =
-\num{1.62e-11}/\num{1.58e-12} = \qty{-10.3}{m/s}$: ke barat.
\end{solution}

\begin{solution}{exo:b3:atmospheric-dynamics:4}
$\beta = 2\Omega\cos\ang{60}/a = \qty{1.14e-11}{m^{-1}.s^{-1}}$;
$\lambda_s = 2\pi\sqrt{10/\num{1.14e-11}} = \qty{5870}{km}$; lingkaran,
\qty{20015}{km}, memuat $n_s = 3.4$ gelombang.
\end{solution}

\begin{solution}{exo:b3:atmospheric-dynamics:5}
$q = f/h$ kekal: di atas pegunungan $\zeta = f(8/10) - f =
\qty{-2.06e-5}{s^{-1}}$, antisiklonik; di atas cekungan $f(12/10) - f =
\qty{+2.06e-5}{s^{-1}}$, siklonik. Tempat kolom mengalami peregangan, di sisi hilir
pegunungan dan di atas cekungan, vortisitas siklonik terbentuk: di sana
depresi cenderung terbentuk.
\end{solution}

\begin{solution}{exo:b3:atmospheric-dynamics:6}
$k = 5/(a\cos\ang{45}) = \qty{1.11e-6}{m^{-1}}$, $\beta/k^2 =
\qty{13.1}{m/s}$: $c = \qty{1.9}{m/s}$, $c_g = \qty{28.1}{m/s}$. Energi
mencakup \qty{5000}{km} dalam \num{2.1}~hari, polanya dalam \num{31}~hari.
\end{solution}

\begin{solution}{exo:b3:atmospheric-dynamics:7}
(a)~$k^2 = \beta/U - l^2 = \num{1.079e-12} - \num{2.74e-13} =
\qty{8.05e-13}{m^{-2}}$, $k = \qty{8.97e-7}{m^{-1}}$: $\lambda =
\qty{7000}{km}$ dan $n = 4.0$, lebih sedikit gelombang dibanding saat $l = 0$. (b)~$l^2 =
\num{1.097e-12} > \beta/U = \num{1.079e-12}$: $k^2$ akan bernilai negatif.
\end{solution}

\begin{solution}{exo:b3:atmospheric-dynamics:8}
$L_d = 0.012\times10^4/\num{1.031e-4} = \qty{1164}{km}$; $\sigma =
0.31\times\num{1.031e-4}\times\num{2e-3}/0.012 = \qty{5.3e-6}{s^{-1}}$, dengan
waktu lipat-$e$ sebesar \num{2.2}~hari; panjang gelombang $2\pi\times1164/1.61 =
\qty{4540}{km}$.
\end{solution}

\begin{solution}{exo:b3:atmospheric-dynamics:9}
$[v^*T^*] = (VT_0/2)\cos\delta$: \qty{25}{K.m/s}, \qty{12.5}{K.m/s} dan 0;
fluks panas \num{25100}, \num{12600} dan \qty{0}{W/m^2}. Ketika
suhu berselisih fase seperempat panjang gelombang dengan angin, sehingga udara hangat
yang bergerak ke selatan sama banyaknya dengan ke utara: tidak ada transpor.
\end{solution}

\begin{solution}{exo:b3:atmospheric-dynamics:10}
Linearisasi di sekitar $U$ dengan $\zeta' = \nabla^2\psi'$ dan $v' =
\partial_x\psi'$: $(\partial_t + U\partial_x)\nabla^2\psi' + \beta\partial_x
\psi' = 0$. With $\psi' \propto \mathrm e^{\mathrm i(kx + ly - \omega t)}$,
$(\omega - Uk)K^2 = -\beta k$, sehingga $\omega = Uk - \beta k/K^2$. Relatif terhadap
udara, kecepatan fasenya adalah $c - U = -\beta/K^2 < 0$ untuk setiap gelombang.
\end{solution}

\begin{solution}{exo:b3:atmospheric-dynamics:11}
$\varphi_H = [5\times9.81\times12\,000\times0.25/(3\times(\num{7.292e-5})^2
\times(\num{6.371e6})^2)]^{1/2} = \qty{0.477}{rad}$, \ang{27.3}.
$u_M = \Omega a\varphi_H^2 = \qty{106}{m/s}$; secara tepat, $\Omega a\sin^2
\varphi_H/\cos\varphi_H = \qty{110}{m/s}$: dua setengah kali
arus jet yang teramati.
\end{solution}

\begin{solution}{exo:b3:atmospheric-dynamics:12}
(a)~Di seluruh pita, antarmuka tergeser sebesar $\eta = sy$, $-L/2 \le y
\le L/2$, dari posisi datarnya. Lapisan dengan kontras massa jenis yang berkurang
tergeser sejauh $\eta$ menyimpan $\tfrac12\rho g'\eta^2$ per satuan luas; dirata-ratakan
di seluruh pita, $\langle\eta^2\rangle = s^2L^2/12$. (b)~$g' = 0.327$:
\num{1.5e5} dan \qty{6.1e5}{J/m^2}, dibandingkan $\tfrac12\times0.7\times15^2
\times10^4 = \qty{7.9e5}{J/m^2}$ energi kinetik: sebanding.
\end{solution}

\begin{solution}{pb:b3:atmospheric-dynamics:1}
\textbf{1.} $f_{45} = \qty{1.031e-4}{s^{-1}}$, $f_{50} =
\qty{1.117e-4}{s^{-1}}$; $\zeta = f_{45} - f_{50} = \qty{-8.6e-6}{s^{-1}}$.

\textbf{2.} $R = V/|\zeta| = 15/\num{8.6e-6} = \qty{1750}{km}$; searah jarum jam,
kembali ke arah selatan.

\textbf{3.} $\zeta = f_{45}(9.5/10) - f_{50} = \qty{-1.38e-5}{s^{-1}}$, lebih
antisiklonik.

\textbf{4.} $\zeta = f_{45} - f_{40} = \qty{+9.4e-6}{s^{-1}}$, siklonik: ia
berbalik ke utara. Bila tergeser ke mana pun, parsel memperoleh vortisitas
yang membalikkannya: efek pemulih, dan osilasi.

\textbf{5.} $\beta = \qty{1.62e-11}{m^{-1}.s^{-1}}$; panjang lingkaran adalah
\qty{28306}{km}.

\textbf{6.} $n = 3$ sampai 8: panjang gelombang 9435, 7076, 5661, 4718, 4044 dan
\qty{3538}{km}; $c = -21.5$, $-5.5$, 1.9, 5.9, 8.3 dan \qty{9.9}{m/s}; $c_g
= 51.5$, 35.5, 28.1, 24.1, 21.7 dan \qty{20.1}{m/s}.

\textbf{7.} $n = 3$ dan 4 bergerak ke barat, $n \ge 5$ ke timur. Gelombang $n = 5$ membutuhkan
$10^6/1.86 = \qty{5.4e5}{s}$, \num{6.2}~hari, untuk bergerak \qty{1000}{km}.

\textbf{8.} Di hilir, ke timur: pada $c_g = \qty{24.1}{m/s}$ energi mencakup
\qty{5000}{km} dalam \num{2.4}~hari.

\textbf{9.} $\lambda_s = 2\pi\sqrt{15/\num{1.62e-11}} = \qty{6050}{km}$,
$n_s = 4.7$.

\textbf{10.} $n_s = 5.7$ untuk \qty{10}{m/s} dan 4.1 untuk \qty{20}{m/s}: angin
baratan yang lebih kuat mendukung gelombang stasioner yang lebih panjang,
lebih sedikit yang melingkari belahan bumi.

\textbf{11.} $n_s = 5.8$ pada \ang{35} LU dan 3.4 pada \ang{55} LU.

\textbf{12.} $\zeta = f(8/10) - f = \qty{-2.06e-5}{s^{-1}}$: pemampatan di atas
pegunungan membuat kolom antisiklonik, sedangkan peregangan di sisi bawah angin
menjadi siklonik. Pegunungan memberlakukan pola vortisitas tetap pada aliran,
yaitu suatu pemaksaan yang tidak bergerak.

\textbf{13.} Pemaksaan tetap menghasilkan respons tunak; respons
paling kuat ketika pola yang dipaksakan sendiri berupa gelombang bebas yang
akan diam di tempat, yaitu resonansi pada $\lambda_s$.

\textbf{14.} $\Lambda = 25/10^4 = \qty{2.5e-3}{s^{-1}}$;
$|\partial T/\partial y| = \Lambda fT/g = \num{2.5e-3}\times\num{1.031e-4}
\times250/9.81 = \qty{6.6}{K}$ per \qty{1000}{km}.

\textbf{15.} $\mathrm{Ri} = 0.011^2/(\num{2.5e-3})^2 = 19$: jauh di atas 0.25,
tidak ada turbulensi geser.

\textbf{16.} $L_d = 0.011\times10^4/\num{1.031e-4} = \qty{1067}{km}$.

\textbf{17.} $2\pi\times1067/1.61 = \qty{4160}{km}$; $28306/4160 = 6.8$
gelombang.

\textbf{18.} $2\pi\times1067/2.40 = \qty{2790}{km}$.

\textbf{19.} $\sigma = 0.31\times\num{1.031e-4}\times\num{2.5e-3}/0.011 =
\qty{7.3e-6}{s^{-1}}$, \num{0.63}~hari$^{-1}$.

\textbf{20.} $\exp(3\times0.627) = 6.6$.

\textbf{21.} Badai, \qty{4160}{km}, lebih pendek daripada
\omterm{prop:b3:atmospheric-dynamics:stationary}{panjang gelombang Rossby stasioner}, \qty{6050}{km}: sebagai gelombang Rossby ia terbawa ke timur
oleh angin baratan, dan badai bergerak ke timur melalui gelombang panjang yang
lebih lambat atau diam.

\textbf{22.} Ke barat seiring ketinggian: hanya dengan begitu udara yang bergerak
ke kutub lebih hangat daripada udara yang bergerak ke khatulistiwa, sehingga panas
bergerak ke kutub dan \omterm{def:b3:atmospheric-dynamics:ape}{energi potensial tersedia} dilepaskan.

\textbf{23.} Gangguan yang tumbuh paling cepat bertambah dengan faktor $e$ setiap
\textbf{\num[text-series-to-math=true]{1.6}~hari} (\num{1.59}~hari menurut
perhitungan).
\end{solution}
